\section{Discussion}
\label{sec:discussion}

\noindent\textbf{Reproduction.} Every baseline we quote was rerun by us and scored with our own metric implementation;
independent PSNR and SSIM implementations agree with the released ones to machine precision. For CIDNet and Retinexformer the reproduced scores agree with the published
ones to \(0.021\) dB except one cell at \(0.081\), a gap we could not attribute
(published/here: CIDNet LOL 23.50/23.498, with perceptual loss 23.8091/23.808, LOL-v2 real 28.1387/28.140 and synthetic 29.5663/29.647 rectified; Retinexformer LOL 27.18/27.169 rectified, Sony 24.44/24.438, LOL-v2 real 27.71/27.689 and synthetic 29.04/29.022 rectified); SNR-Net, LLFormer and URetinex-Net in its
reference-ratio setting reproduce their published LOL scores to \(0.003\) dB, the two diffusion models to \(0.13\) dB (GSAD, under rectification, beyond its single-run
spread) and \(0.34\) dB (LightenDiffusion, within its seed spread), and
Zero-DCE++ has no published LOL score; this is what licenses reading the decomposition as a
property of the method.



\noindent\textbf{Three protocol choices move the score by as much as the margins do.} First,
reference-mean rectification sets the output's brightness from the reference. On the Sony benchmark it is worth \(1.16\) dB, the whole-frame special case of the
spatially varying term priced in Sec.~\ref{sec:results}, and for GSAD on LOL \nGSADrectGain{} dB; a related setting, URetinex-Net's
reference-derived ratio, is worth
\nURratioGain{} dB (\nURpub{} published against \nURfixed{} at its fixed ratio, Table~\ref{tab:decomp}). Retinexformer's documentation recommends against the setting while CIDNet reports its
headline number with it, so the two are not comparable; the latter's script further scales LOL-v2 outputs by a per-checkpoint constant chosen
against the reference (about three decibels unrectified). All three are reasons to rerun baselines, not criticisms of a method. Second, the 8-bit conversion rule: one pipeline saves outputs by truncation and the other by rounding, confirmed at the pixel
level on saved files. Switching the rule alone is worth between \nBITlo{} and \nBIThi{} dB on the three LOL
checkpoints. Third, the dataset variant: the Sony extreme-dark set is distributed in at least two renditions whose pixel conventions
differ~\cite{retinexformer,cidnet,cidnetplus},
and the checkpoint of~\cite{cidnet} scores \nCROSSpsnr{} dB on the other ($598$ frames, rounding) against the \(22.904\)
reported for it on its own.

\noindent\textbf{What the decomposition implies for method design.} The headroom on Sony is \nHRglob{} dB for a global gain and \nHRchan{} dB
more for a per-channel gain, \nHRblk{} dB for a per-block gain and \nHRaff{} dB for a
per-block affine map at $16$ pixels over the whole-frame fit (both control-corrected),
\nHRbandls{} dB for per-band amplitude, and about \nHRnoise{} dB for the noise share estimated under independence. The rule that follows is to spend effort on the component with the largest
headroom, testing its predictors before building on it: on Sony the spatially varying low-frequency term, on LOL, by share of squared error, the global term for every method but Retinexformer
(Table~\ref{tab:decomp}).
Single-frame high-frequency restoration is worth a few percent under the noise-corrected
ceiling, and forcing amplitude up makes it worse. The low-frequency term resists the predictors we tried; the multi-frame margin exists in the input but this single-frame model cannot take it,
so taking it would require training on burst statistics. A testable prescription: train on $k$-frame means and test whether the $k=8$ low-frequency energy drops below the single-frame value.

\noindent\textbf{Limitations.} All quantities are measured on the 8-bit sRGB rendition the benchmarks distribute, so every recoverability statement is about that rendition, not sensor data. Quantisation to 8 bits, and sensor noise that is not additive~\cite{eld2020}, depart from
the additive uncorrelated-noise model behind the noise correction; the check of Sec.~\ref{sec:results} fails in the lowest band at $0.04$ and $0.1$\,s, and
cells with $N_j/P_j>0.9$ are reported as unstable. The identifiability test is a
distortion-axis statement about conditional-mean regressors~\cite{blau2018}, so a perceptually
optimised method can fall below the input on $\rho_j$ without being wrong.
The reference's noise floor was not measured and is treated as zero.
The band-wise test covers Retinexformer on Sony and CIDNet and Retinexformer on LOL,
while Table~\ref{tab:decomp} decomposes every method; CIDNet's Sony residual was not decomposed because the variant it was trained on, in a
different pixel convention, is not the one we could download. The noise-corrected ceiling exists only on Sony, which has repeated captures; band
correlations are global per frame, so a band can be labelled exhausted while local
disagreement remains; the multi-frame result uses static scenes only.

\section{Conclusion}
\label{sec:conclusion}

We decomposed the residual of reproduced enhancement models into components with oracle
ceilings, control-corrected where spatial, and tested band by band whether the error is
discarded information or information the rendition never carried. On both benchmarks the residual after the global gain sits at low frequency, a spatially varying correction of it is worth more than a decibel, and the output is closer to a blurred reference than to the reference. That term resists the
predictors we tried, this single-frame checkpoint cannot consume a burst, and code and
numbers are released.
